\begin{center}
    \section*{\fontsize{20}{20}\selectfont Chapter 2}
\end{center}
\vspace{10mm}
\section{Background Study}

Autism Spectrum Disorder (ASD) is a neurodevelopmental disorder and it puts an affect to an children's ability to interact and communicate with others and also exhibits repetitive behaviors. Early detection of behaviors associated with autism is essential for successful intervention and developmental cure and support in the long run. Unfortunately, most current approaches in our country to diagnose autism are highly dependent on expert observation and manual behavioral analysis.

With the rising of new technologies and methods in artificial intelligence, interested researchers and scientists in the respected field have been using machine learning and computer vision approaches for the detection of autistic patterns and analysis of the related behaviors. Researchers have been showing great interest in using pose estimation, video analysis for different repetitive action recognition and classification approaches for the analysis of behaviors objectively. This chapter presents a review of the existing literature related to the detection of autism, analysis of behaviors using pose estimation, and application of machine learning approaches for the analysis of these behaviors which are going to be used for the proposed system.


\subsection{Literature Review}

There are many research papers that have investigated the usage of machine learning and computer vision techniques for applications related to behavioral analysis. 

In the research paper by Rajpurkar et al. \cite{rajpurkar2022ai}, the authors discussed the increasing involvement of artificial intelligence with regard to medical decision support systems. They reduced the manual workload and improved the validations consistently which supported the healthcare professionals and their system was AI based. Hence, the idea is very relevant to our work as our thesis is also about specific behavioral screening of an autsim child via Aritificial Intelligents footage analysis.

Duda et al. \cite{duda2016use} studied the use of machine learning classifiers in screening autism using behavioral symptoms and promising results were gained in screening of ASD in its early stages. Although their metrice on perfomance was through reducing behavioral questions, they also maintained good classification performance. This is relevant to our thesis idea as we selected repetitive behavioral patterns as a strong indicator for autism suspicion in early stages.

Another study by Thabtah \cite{thabtah2018machine} surveyed various machine learning approaches in diagnosing autism and identified some of the challenges in handling imbalanced datasets and feature selecting. Their paper clearly represented the importance of data preprocessing and how much influencial it is to improving autism related predictions. They did it by using sequence normalization and augmentation in proper methods which played a clear role as a reference for our thesis. 

Pose-based human action classification has also been greatly explored. Cao et al. \cite{cao2017realtime} proposed a real-time multi-person pose estimating framework that had a large effect on subsequent behavior analysis systems. 

Liu et al. \cite{liu2019skeleton} employed features arising from skeletal movements for different action recognition and also demonstrats that pose landmarks can be very useful for representing human motions. They also have an pose estimation framework which later became as one of the foundational works for extracting body landmarks from video footages like ours. This is closely similar our work of MediaPipe pose usage to collect the skeleton basis movements featured from children videos.

In the context of autism studies, Hashemi et al. \cite{hashemi2014computer} used computer vision to analyze the patterns in the facial and bodily movement of children having autism. This paper has opened a great door of opportunity for us by referencing that subtle non verbal behavioral cues which can be hidden to see by naked eye can be captured with computer vision. This worked as a super insight suggestion for us and motivated us for pose based approaching of our research that detects repetitive autism patterns from short video clips.


Voss et al. \cite{voss2019emotion} studied emotion recognition in children with autism by using video data, emphasizing on the role of temporal modeling. This paper helped us to keep a close attention to the temporal dependency learnt from the sequences extracted from the videos and it it highly important for autism related behavioural pattern analysis and its supported the evaluation our deep learning models like CNN-BiLSTM which relies on temporal sequence learning for repeated motion analysis.

As for Deep learning techniques ,Ord\'{o}\~{n}ez et al. \cite{ordonez2016deep} have successfully utilized LSTM networks for activity recognition in respected time series data. Their worked provided evidence that deep learning model which are sequential can learn motion dependencies over the times and this concept heavily influenced our deeplearning AI architechtures like CNN-BiLSTM.

More recently, Bi-directional LSTM networks have been utilized for learning long-term temporal dependencies from human behavior sequences \cite{zhao2020bidirectional}. And most importantly their bidirectional learning process aided us to understand the both previous and future frame relationships. This is really important for our research because autism specific behavioral traits contains repetitive forward and backward moving patterns.

Explainable artificial intelligence increases the trust and reliance in machine learning models. Lundberg and Lee \cite{lundberg2017unified} proposed SHAP to interpret the model predictions. The proposed method was then used in healthcare-related classification problems. The framworks shows and identifies which body landmarks contributed the most to the final prediction which is highly influencial for our thesis because it helps us to take our final result on which body landmark and pose movements were most dominant in the classfication.


Ribeiro et al. \cite{ribeiro2016should} highlighted the importance of model interpretability, especially in healthcare-related problems. Their work binds the tie and trust among the healtcare professional and researchers highlighting interpretable AI and it is indeed the reason why we included explainability in our final model analysis.

In developing countries, automated systems are very important and feasible. Hossain et al. \cite{hossain2023smart} proposed a smart decision support system more fouced on the integration of machine learning for real-world applications. Their work showed the adaptation of intelligent systems as a practical deployment in the developing regions. This is crucial intention of our thesis as we also intend to screen and support the austism affected children early for detection in resource limited environment like ours developing region.

The above studies altogether is really worried about the importance of a pose-based and data-driven approach for autism behavior classification.


\begin{table}[h]
\centering
\caption{Comparison of Related Works}
\begin{tabular}{|p{3cm}|p{3cm}|p{3cm}|p{3cm}|}
\hline
\textbf{Author (Year)} & \textbf{Data Type} & \textbf{Method} & \textbf{Limitation} \\ \hline
Duda et al. (2016) \cite{duda2016use} & Behavioral data & ML classifiers & Limited features \\ \hline
Thabtah (2018) \cite{thabtah2018machine} & Survey data & ML review & No video analysis \\ \hline
Hashemi et al. (2014) \cite{hashemi2014computer} & Video & CV techniques & Small dataset \\ \hline
Cao et al. (2017) \cite{cao2017realtime} & Video & Pose estimation & Not ASD-focused \\ \hline
Liu et al. (2019) \cite{liu2019skeleton} & Skeleton data & Action recognition & General actions \\ \hline
Ord\'{o}\~{n}ez et al. (2016) \cite{ordonez2016deep} & Time-series & LSTM & No pose data \\ \hline
Zhao et al. (2020) \cite{zhao2020bidirectional} & Video & BiLSTM & High complexity \\ \hline
Voss et al. (2019) \cite{voss2019emotion} & Video & Emotion recognition & Limited behaviors \\ \hline
Lundberg et al. (2017) \cite{lundberg2017unified} & Tabular & Explainable AI & Not behavioral \\ \hline
Ribeiro et al. (2016) \cite{ribeiro2016should} & ML models & Model explanation & No temporal data \\ \hline
Rajpurkar et al. (2022) \cite{rajpurkar2022ai} & Healthcare data & AI systems & Conceptual \\ \hline
Hossain et al. (2023) \cite{hossain2023smart} & IoT + ML & Decision systems & Not ASD-specific \\ \hline
Sadman et al. (2023) \cite{sadman2023password} & Security data & ML framework & Different domain \\ \hline
\end{tabular}
\end{table}


\subsection{Problem Analysis}

so by reviewing the literature , it is clear that although there have been successful applications of machine learning and computer vision in field of healthcare and behavior analysis, there have been limitations in autism research. Most studies have been static in nature and have failed to fully account for temporal movement patterns. Video-based studies have been characterized by small databases, lack of generalization and poor interpretability.

Moreover, the available methods hardly employ any the integration of pose-based feature extraction with machine learning explainability. The absence of such mechanisms in decision making further limits the usability of the methods in real-world applications. Such gaps have emphasized the need for developing an integrated system that is capable of analyzing temporal behavioral patterns while offering reliability in classification with interpretability.


From the background study, it is clear that there is increasing interest in the use of machine learning and computer vision for the analysis of autism behavior. Even though there is already research on the feasibility of the approach, there have been also limitations in the approach due to issues such as temporal modeling, explainability and dataset limitations. This research aims to bridge that gap by introducing a new approach to autism behavior classification using a pose-based machine learning approaching.
